\documentclass[11pt]{letter}
\usepackage[margin=1in]{geometry}
\usepackage{mathptmx}
\usepackage{amsmath}
\usepackage{hyperref}
\signature{Ehsan Roohi\\Corresponding author\\Department of Mechanical and Industrial Engineering\\University of Massachusetts Amherst}
\address{Department of Mechanical and Industrial Engineering\\University of Massachusetts Amherst\\Amherst, MA 01003, USA}
\date{August 30, 2026}

\begin{document}
\begin{letter}{Editor\\\emph{Physics of Fluids}}
\opening{Dear Editor,}

Please consider our substantially revised manuscript, ``Bulk-to-Wall Observability in Rarefied Hypersonic Flow over Triangular Protrusions: Full-Range Moment Limits and Incident Half-Range Sufficiency,'' by Elyas Lekzian and Ehsan Roohi, for publication in \emph{Physics of Fluids}.

The earlier version centered on a model-conditional spatial-support diagnostic, $R_{95}$, extracted from DSMC macroscopic fields. The review correctly identified that this evidence was useful but not sufficient for a general physical claim about wall-load observability, particularly for shear. We have therefore retained the full validation and reproducibility structure while rebuilding the scientific center of the manuscript around new DSMC data and a kinetic sufficiency test.

The revised paper now makes three main contributions. First, an equal-capacity blind interpolation experiment compares primitive variables, the complete momentum-flux tensor, and heat-flux augmentation. The stress tensor improves pressure relative to primitive fields, but finite full-range states do not improve transferable signed shear. Second, a constructive positive-distribution example proves that full-range moments through degree three can be identical while incident half-range pressure differs and half-range shear changes sign. Third, SPARTA collision tallies show that the incoming half-range state, combined with the known diffuse-wall kernel, reconstructs co-temporal DSMC pressure to 0.15--0.28\% NRMSE and signed shear to 1.30--7.35\% across six geometry--rarefaction cases. Independent 40,000-step windows retain 1.68--2.85\% pressure error and identify the forward-facing shear cases as sampling-noise limited rather than concealing them behind absolute shear.

We believe the importance of this result is broader than one surrogate architecture. A kinetic--continuum interface or learned wall closure can fail even when it receives accurate low-order macroscopic fields, because full-range moments combine molecules approaching and leaving the wall. The incident half-range result provides a physically interpretable upper-bound state against which practical directional-moment or reduced-distribution closures can be designed. The neural surrogate remains in the paper as a validated emulator and controlled comparison tool; it no longer carries the central kinetic claim. Likewise, $R_{95}$ is retained only as a secondary model-conditional diagnostic.

The manuscript is a strong fit for \emph{Physics of Fluids}: it combines rarefied-gas dynamics, DSMC, gas--surface interaction, uncertainty-aware surface loading, and data-driven closure assessment. All figures have been rebuilt in publication-quality landscape form. The reproducibility package contains case manifests, SPARTA inputs, block-resolved moment and collision outputs, exact commands, packed arrays, numerical quality-control reports, source scripts, vector figures, and checksums.

This manuscript is original, is not under consideration elsewhere, and has been approved by both authors. The authors declare no conflict of interest and no external funding for this work.

Thank you for your consideration.

\closing{Sincerely,}
\end{letter}
\end{document}
